\documentclass[11pt]{article}

\usepackage[margin=1in]{geometry}
\usepackage{amsmath, amssymb, amsthm, mathtools}
\usepackage{bm}
\usepackage{enumitem}

\title{Recovering Row-Wise Categorical Vectors from Sampled $F$-Matrices}
\author{}
\date{July 21, 2026}

\begin{document}

\maketitle

\section{Goal}

We consider the following simplified problem. Let $n = 10$, where $n$ denotes the
dimension of the lower-triangular $F$-matrix. Suppose there exist true row-wise
categorical vectors
\[
\pi_i = (\pi_{i1}, \ldots, \pi_{i,i-1}), \qquad i = 3, \ldots, 10,
\]
with
\[
\pi_{ij} \ge 0,
\qquad
\sum_{j=1}^{i-1} \pi_{ij} = 1.
\]
Using these vectors, we generate $M = 10{,}000$ independent $F$-matrices
\[
F^{(1)}, \ldots, F^{(M)}.
\]
The objective is to recover the unknown categorical vectors
\[
\{\pi_3, \pi_4, \ldots, \pi_{10}\}
\]
from the sampled $F$-matrices.

\section{Generative Construction}

The matrix $F$ is built row by row. The first two rows are fixed:
\[
F_{11} = 2,
\qquad
F_{21} = 1,
\qquad
F_{22} = 3,
\]
and the diagonal is always fixed as
\[
F_{ii} = i+1, \qquad i = 1, \ldots, n.
\]

For each row $i \in \{3, \ldots, n\}$, let
\[
r_{i-1} = (F_{i-1,1}, \ldots, F_{i-1,i-1})
\]
denote the previous row. Define the row increments
\[
\Delta_1(r_{i-1}) = F_{i-1,1},
\qquad
\Delta_j(r_{i-1}) = F_{i-1,j} - F_{i-1,j-1},
\qquad j = 2, \ldots, i-1.
\]
The valid breakpoint set for row $i$ is
\[
V_i(r_{i-1})
=
\{j \in \{1, \ldots, i-1\} : \Delta_j(r_{i-1}) > 0\}.
\]

Intuitively, a breakpoint is valid exactly when that position corresponds to an
upward step in the previous row. This removes both leading zeros and flat spots.

Conditional on the previous row, the breakpoint $S_i$ is sampled from the masked
and renormalized categorical distribution
\[
\mathbb{P}(S_i = j \mid r_{i-1})
=
\frac{\pi_{ij}}{\sum_{k \in V_i(r_{i-1})} \pi_{ik}},
\qquad j \in V_i(r_{i-1}),
\]
and $\mathbb{P}(S_i = j \mid r_{i-1}) = 0$ for $j \notin V_i(r_{i-1})$.

Once $S_i$ is sampled, row $i$ is obtained deterministically from row $i-1$ via
a suffix drop:
\[
F_{ij} =
\begin{cases}
F_{i-1,j}, & j < S_i, \\
F_{i-1,j} - 1, & S_i \le j \le i-1,
\end{cases}
\qquad
F_{ii} = i+1.
\]

Thus each row is governed by a single latent breakpoint rather than by independent
cell-wise Bernoulli variables.

\section{Observed Data and Latent Breakpoints}

Suppose we observe the sampled matrices
\[
F^{(1)}, \ldots, F^{(M)}.
\]
For each sample $m$ and row $i$, define the row difference vector
\[
d_i^{(m)}
=
\bigl(
F_{i-1,1}^{(m)} - F_{i,1}^{(m)},
\ldots,
F_{i-1,i-1}^{(m)} - F_{i,i-1}^{(m)}
\bigr).
\]

Under the suffix-drop construction, this vector must have the form
\[
d_i^{(m)} = (0, \ldots, 0, 1, \ldots, 1).
\]
Therefore the latent breakpoint is directly recoverable from the observed matrix:
\[
S_i^{(m)} = \min\{j : d_{i,j}^{(m)} = 1\}.
\]
At the same time, the valid breakpoint set is also recoverable from the previous
row:
\[
V_i^{(m)}
=
\left\{
j : \Delta_j\bigl(r_{i-1}^{(m)}\bigr) > 0
\right\}.
\]

Hence for each row $i$ and each sample $m$, the observed $F$-matrix determines
both
\[
S_i^{(m)}
\qquad \text{and} \qquad
V_i^{(m)}.
\]

\section{Likelihood for Recovering the Categorical Vectors}

Because the breakpoints are recoverable from the observed matrices, this problem
is not yet a variational inference problem. It is a masked categorical
maximum-likelihood problem.

For a fixed row $i$, the likelihood contribution of sample $m$ is
\[
\mathbb{P}(S_i^{(m)} \mid V_i^{(m)}, \pi_i)
=
\frac{\pi_{i,S_i^{(m)}}}
{\sum_{k \in V_i^{(m)}} \pi_{ik}}.
\]
Therefore the row-wise log-likelihood is
\[
\ell_i(\pi_i)
=
\sum_{m=1}^M
\left[
\log \pi_{i,S_i^{(m)}}
-
\log\left(
\sum_{k \in V_i^{(m)}} \pi_{ik}
\right)
\right].
\]

The full log-likelihood is the sum across rows:
\[
\ell(\pi_3, \ldots, \pi_{10})

=
\sum_{i=3}^{10} \ell_i(\pi_i).
\]
This separates by row, so each $\pi_i$ can be estimated independently.

\section{Unconstrained Parameterization}

A convenient parameterization uses logits $\theta_i = (\theta_{i1}, \ldots,
\theta_{i,i-1})$ with
\[
\pi_{ij}
=
\frac{\exp(\theta_{ij})}
{\sum_{\ell=1}^{i-1} \exp(\theta_{i\ell})}.
\]
Substituting this into the row-wise log-likelihood gives
\[
\ell_i(\theta_i)
=
\sum_{m=1}^M
\left[
\theta_{i,S_i^{(m)}}
-
\log\left(
\sum_{k \in V_i^{(m)}} \exp(\theta_{ik})
\right)
\right].
\]

The gradient with respect to $\theta_{ij}$ is
\[
\frac{\partial \ell_i}{\partial \theta_{ij}}
=
\sum_{m=1}^M
\left[
\mathbf{1}\{S_i^{(m)} = j\}
-
\mathbf{1}\{j \in V_i^{(m)}\}
\frac{\exp(\theta_{ij})}
{\sum_{k \in V_i^{(m)}} \exp(\theta_{ik})}
\right].
\]
This makes gradient ascent, Adam, or any other smooth optimizer straightforward.

\section{Estimation Algorithm}

For the present setup with $n=10$ and $M=10{,}000$ sampled matrices, the
estimation problem can be written as follows.

\begin{enumerate}[leftmargin=2em]
    \item For each sampled matrix $F^{(m)}$ and for each row $i = 3, \ldots, 10$:
    \begin{enumerate}[label=(\alph*), leftmargin=2em]
        \item compute the breakpoint $S_i^{(m)}$ from the first location where
        row $i$ differs from row $i-1$ by $1$;
        \item compute the valid set $V_i^{(m)}$ from the positive increments of
        row $i-1$.
    \end{enumerate}
    \item For each row $i = 3, \ldots, 10$, maximize
    \[
    \ell_i(\theta_i)
    =
    \sum_{m=1}^M
    \left[
    \theta_{i,S_i^{(m)}}
    -
    \log\left(
    \sum_{k \in V_i^{(m)}} \exp(\theta_{ik})
    \right)
    \right].
    \]
    \item Convert the fitted logits to probabilities using the softmax map:
    \[
    \widehat{\pi}_{ij}
    =
    \frac{\exp(\widehat{\theta}_{ij})}
    {\sum_{\ell=1}^{i-1} \exp(\widehat{\theta}_{i\ell})}.
    \]
    \item Validate the fitted vectors by simulating new $F$-matrices from
    $\widehat{\pi}_3, \ldots, \widehat{\pi}_{10}$ and comparing the induced
    breakpoint frequencies and matrix summaries to the original sample.
\end{enumerate}

\section{Remarks}

\begin{itemize}[leftmargin=2em]
    \item For $n=10$, the row-wise categorical vectors have lengths
    $2,3,\ldots,9$. This means there are $44$ raw logits but only
    $36$ free probability parameters after accounting for the sum-to-one
    constraint in each row.
    \item The masking is essential. Naive empirical frequencies of the observed
    breakpoints are generally biased because some breakpoints are unavailable in
    some rows.
    \item This estimation problem is a useful first stage before moving to a
    full variational inference setup in which the $F$-matrices themselves are not
    observed.
\end{itemize}

\end{document}
